\begin{theorem}[Proved: the functor equivalences with their exact meaning]\label{cat:I12}
For accessible categories $E,F$, the following are equivalent.
\begin{enumerate}
\item $T:E\to F$ is accessible.
\item At some common accessibility rank $\lambda$, it is the Ind
extension of a functor $E_\lambda\to F_\lambda$.
\item It is the pointwise left Kan extension of its restriction to a
small full subcategory of $E$.
\end{enumerate}
In (3), ``pointwise'' means that the canonical colimits over $C/X$
exist and compute the extension; this specifies the meaning of the
manuscript's $i_!^F$ on scan 055.
Moreover, if $E$ is $\kappa$-accessible and a locally small category
$F$ has $\kappa$-filtered colimits, restriction is an equivalence
\[
 \operatorname{Fun}_{\kappa\text{-filt}}(E,F)
   \xrightarrow{\sim}\operatorname{Fun}(E_\kappa,F).
\]
Here the left side consists of functors preserving those colimits,
with all natural transformations.  No presentability condition on
the values in $F$ is needed for this last equivalence.
\end{theorem}
\begin{proof}
Choose a small skeleton of $E_\kappa$.  The extension of a functor
$h:E_\kappa\to F$ is
\[
 \overline h(X)=\operatorname{colim}_{(c\to X)\in E_\kappa/X}h(c).
\]
The indexing category is $\kappa$-filtered by
Theorem~\ref{cat:I1}, so the colimit exists.  Its universal property
identifies maps $\overline h(X)\to Z$ with natural transformations
\[
 \operatorname{Hom}_E(-,X)|_{E_\kappa}
       \longrightarrow \operatorname{Hom}_F(h(-),Z).
\]
For a $\kappa$-filtered colimit of $X$'s the presheaf on the left is
their pointwise colimit.  Mapping this colimit to the right-hand
presheaf gives the limit of the corresponding sets of natural
transformations.  Thus $\overline h$ preserves the required colimits.
At a presentable object $c$ the slice has terminal object $1_c$, so
the restriction is $h$.  Conversely any preserving functor is
computed by these canonical presentations.  The same colimit
property extends every natural transformation uniquely.  This proves
the displayed equivalence and the pointwise Kan assertion at this rank.

For an accessible $T$, first choose a cardinal $\nu$ above given
accessibility ranks for $E,F$ and a preservation rank for $T$.
Theorem~\ref{cat:I11}, applied separately to $E$ and $F$, makes
$(2^\nu)^+$ a common accessibility rank, and $T$ preserves filtered
colimits at that rank because it is larger than its preservation rank.
Theorem~\ref{cat:I8} now supplies a common rank
at which $T$ also preserves presentables; Proposition~\ref{cat:I5}
then gives (2).  Condition (2) implies (1) and (3) by the preceding
paragraphs.
Suppose (3) holds for a small full category $C$.  Each object of an
accessible category is presentable at some rank: write it as a small
filtered colimit of presentables and take a regular cardinal larger
than the indexing category's arrow set.  The Hom-limit calculation
used in Theorem~\ref{cat:I11} proves this assertion directly.
Since $C$ is small, choose one regular $\lambda$ making all its
objects $\lambda$-presentable and such that $E,F$ admit
$\lambda$-filtered colimits.  The same natural-transformation formula,
now with $C$ in place of $E_\kappa$, shows that the pointwise Kan
extension preserves those colimits.  This is (1).
\end{proof}

